%% ============================================================================
%% paper.tex
%% "The WKB Geometry of Newton Polygons: A Functorial Classification Theory"
%% Version: v1.0  (camera-ready candidate for Zenodo upload v1.0)
%% Companion dataset: doi:10.5281/zenodo.XXXXXXX  (PLACEHOLDER)
%% ============================================================================

\documentclass[11pt,a4paper]{article}

%% --- Geometry & typography -------------------------------------------------
\usepackage[margin=1in]{geometry}
\usepackage[T1]{fontenc}
\usepackage[utf8]{inputenc}
\usepackage{lmodern}
\usepackage{microtype}

%% --- Math ------------------------------------------------------------------
\usepackage{amsmath, amssymb, amsthm, mathtools}
\usepackage{mathrsfs}
\usepackage{stmaryrd}

%% --- Tables, lists, graphics, color ---------------------------------------
\usepackage{booktabs}
\usepackage{array}
\usepackage{enumitem}
\usepackage{graphicx}
\usepackage{xcolor}

%% --- Hyperlinks & cross-refs (load late) ----------------------------------
\usepackage[colorlinks=true,
            linkcolor=blue!50!black,
            citecolor=blue!50!black,
            urlcolor=blue!50!black]{hyperref}
\usepackage[capitalize,noabbrev]{cleveref}

%% --- Theorem environments --------------------------------------------------
\theoremstyle{plain}
\newtheorem{theorem}{Theorem}[section]
\newtheorem{lemma}[theorem]{Lemma}
\newtheorem{proposition}[theorem]{Proposition}
\newtheorem{corollary}[theorem]{Corollary}

\theoremstyle{definition}
\newtheorem{definition}[theorem]{Definition}
\newtheorem{example}[theorem]{Example}
\newtheorem{construction}[theorem]{Construction}

\theoremstyle{remark}
\newtheorem{remark}[theorem]{Remark}
\newtheorem{notation}[theorem]{Notation}

%% --- Macros (notation lock) -----------------------------------------------
\newcommand{\op}{\mathop{\square}}                % abstract irregular operator
\newcommand{\opD}{\mathrm{D}}                     % d/dx
\newcommand{\opS}{\mathrm{S}}                     % shift
\newcommand{\opT}{\mathrm{T}}                     % q-shift
\newcommand{\LCinf}{\mathrm{LC}_\infty}
\newcommand{\Roots}{\mathrm{Roots}}
\newcommand{\NL}{N(L)}
\newcommand{\NLa}{N(L)}
\newcommand{\Cl}{\mathcal{C}}
\newcommand{\Fclass}{\mathsf{F}}
\newcommand{\MW}{\mathrm{MW}}
\newcommand{\Sys}{\mathrm{Sys}}
\newcommand{\Sym}{\mathrm{Sym}}
\newcommand{\Symm}{\mathfrak{S}}
\newcommand{\disc}{\mathrm{disc}}
\newcommand{\Aset}{\mathcal{A}}
\newcommand{\Acov}{\mathcal{A}^{\mathrm{cov}}}
\newcommand{\Abase}{\mathcal{A}^{\mathrm{base}}}
\newcommand{\Dij}{\Delta_{ij}}
\newcommand{\ie}{i.e.\ }
\newcommand{\eg}{e.g.\ }
\newcommand{\cf}{cf.\ }
\newcommand{\Zlc}{\mathbb{Z}[\zeta_q]}
\newcommand{\Qlc}{\mathbb{Q}(\zeta_q)}
\newcommand{\sloped}[2]{-#1/#2}
\DeclareMathOperator{\arginline}{arg}
\newcommand{\Argm}{\arginline}

%% --- Title & authors -------------------------------------------------------
\title{The WKB Geometry of Newton Polygons:\\
       A Functorial Classification Theory}

\author{%
  Papanokechi\thanks{Independent researcher. ORCID:
  \href{https://orcid.org/0009-0000-6192-8273}{\texttt{0009-0000-6192-8273}}.
  Code and data repository:
  \url{https://github.com/papanokechi}.}
}

\date{\today\hspace{1ex}\textsc{(preprint v1.0)}}


\begin{document}
\maketitle

%% ============================================================================
%% ABSTRACT
%% ============================================================================
\begin{abstract}
Local linear operators with one irregular singularity---differential,
difference, or $q$-difference---are organised, via the WKB ansatz over the
dominant Newton edge, by a discrete invariant: the \emph{exponent
constellation} $\Cl(L) \subset \mathbb{C}$. We promote this organisation to
a functorial classification. Writing
$\Fclass \colon L \mapsto \Cl(L)$ from scalar local models with one dominant
Newton edge to finite multisets in $\mathbb{C}$ taken up to global rotation,
we determine the image of $\Fclass$ and the structure of its fibres.

The main structural law is algebraic. If the dominant Newton edge has slope
$-p/q$ in lowest terms and the edge contains $r{+}1$ lattice points, then
the WKB characteristic polynomial factors as
$\chi(c) = \chi_w(c^q)$ with $\chi_w \in \mathbb{C}[w]$ of degree $r$.
Consequently
$\Cl(L) = \bigsqcup_{w_i \in \Roots(\chi_w)} \mu_q \cdot w_i^{1/q}$,
a $\mu_q$-equivariant pullback of the inner spectrum along $c \mapsto c^q$.
The image of $\Fclass$ is a stratified space whose generic stratum, indexed
by $(p, q, r)$, contains finite-multiset configurations of $r$ concentric
$q$-gons; deeper strata correspond to equal-modulus collisions, multiplicities
in $\chi_w$, and zero roots.

An automated sweep over roughly $96$ operators---five slopes, three ranks,
four valuation patterns, three equation types, plus stress-test probes and
multi-edge / system / nested-ramification extensions---yields \textbf{$23$}
distinct constellation classes. Beyond the single-edge case we prove:
(M1) a per-edge factorisation along the slope filtration; (N1) an iterated
$\mu_q$-pullback law for nested inner polynomials; (D2) a prime-$q$ Kummer
dichotomy distinguishing algebraically independent and cyclotomic strata;
and (S1) a block-diagonal pullback for systems. The discrete classification
admits an analytic refinement: (AT1) anti-Stokes directions are determined
by $\arg \Dij$ on the ramified cover; (AT3) nested ramification induces an
extra $\mu_{q'}$-action on wild character variety strata while the minimal
Levelt--Turrittin cover stays at $q$; (AT5) multi-edge Stokes data are
slope-graded; (AT7) the slope numerator $p$, invisible to $\Cl(L)$, is
visible to the analytic anti-Stokes skeleton.

The contribution is a classification of the functor itself, not of the
operators feeding it: a discrete geometry of WKB images, with explicit
generic and special strata, intrinsic monodromy, and equation-type-invariant
combinatorics.
\end{abstract}

\bigskip
\noindent\textsc{Keywords:} WKB analysis; Newton polygon; irregular
singularity; formal classification; Levelt--Turrittin; Stokes phenomenon;
anti-Stokes directions; wild character variety; functorial classification;
$\mu_q$-pullback.

\bigskip
\noindent\textsc{MSC2020:} 34M40 (primary); 34M30, 34M35, 39A06, 39A13
(secondary).

\tableofcontents

\bigskip\hrule\bigskip


%% ============================================================================
%% §1  INTRODUCTION
%% ============================================================================
\section{Introduction}\label{sec:intro}

\subsection{What we classify, and why a functor}\label{ssec:intro-functor}

Linear operators with an irregular singularity occupy a peculiar middle
ground in the analytic theory of differential and difference equations. They
are simple enough that a great deal of structure---Newton polygons, formal
solutions, Stokes matrices---can be read off finitely many coefficients,
yet rich enough that no single invariant has proved adequate to organise
the resulting zoo. The present paper proposes to study not individual
operators but the \emph{map} that extracts the discrete WKB content of an
irregular singularity, and to classify its image.

Concretely, attach to a scalar linear operator $L$ with one irregular
singularity at $x = \infty$ its dominant Newton edge of slope $-p/q$
($\gcd(p, q) = 1$), apply the WKB ansatz $f \sim \exp(\int c\, x^{p/q-1}\,
\mathrm{d}x)$ along that edge, and extract the leading balance: the
\emph{exponent constellation}
$\Cl(L) \subset \mathbb{C}$, a finite multiset. This assignment
$\Fclass \colon L \mapsto \Cl(L)$ is the central object of the paper, and
we call it the \emph{constellation functor}. It is invariant of equation
type ($\opD$, $\opS$, or $\opT$) at the combinatorial level; the analytic
meaning of an element of $\Cl(L)$ depends on the equation type, but the
multiset does not.

\subsection{The structural law}\label{ssec:intro-law}

The core observation is that $\Cl(L)$ is not an arbitrary multiset.
\Cref{thm:main} below establishes the algebraic identity
\[
   \chi(c) \;=\; \chi_w(c^q),
\]
where $\chi$ is the WKB characteristic polynomial built from the leading
coefficients along the dominant edge and $\chi_w$ is the \emph{inner
polynomial}: a degree-$r$ polynomial whose coefficients are exactly the
$\LCinf$ data of the edge in lowest-degree-first order. Solving $\chi = 0$
amounts to solving $\chi_w$ and then pulling back the roots along
$c \mapsto c^q$; the upshot is that $\Cl(L)$ decomposes canonically as a
$\mu_q$-equivariant union of $r$ orbits.

This factorisation is at once an obstruction (no constellation outside the
image of $\Fclass$ has it) and a recipe (every constellation in the image
is so built). Section~\ref{sec:main} states the law as a theorem with two
corollaries (ring decomposition; Vieta monodromy) and a remark on the
degeneracy strata.

\subsection{From a law to a classification}\label{ssec:intro-classification}

The structural law turns the classification problem into a combinatorial
one. Up to global rotation, a constellation is recovered from $(p, q, r)$
and the modulus-with-multiplicity partition $\pi$ of the roots of $\chi_w$.
A computational sweep, recorded in the dataset companion to this paper
\cite{wkb_dataset_v1}, enumerates these classes for
$p/q \in \{1/2, 1/3, 1/4, 2/3, 3/2\}$ and rank $r \in \{1, 2, 3\}$ across
three equation types and a battery of stress-test probes: there are exactly
$23$ classes. These populate the open stratum and four named degenerate
strata (equal modulus; resonant multiplicity; zero root; accidental
symmetry); \Cref{sec:classes} catalogues them.

The classification is then extended in several directions:
\begin{itemize}[leftmargin=*]
\item beyond the single edge: per-edge factorisation along the slope
      filtration (\Cref{thm:M1}, \S\ref{sec:multiedge});
\item beyond a single $\mu_q$-orbit pulling: nested ramification with
      iterated pullbacks (\Cref{thm:N1}, \S\ref{sec:multiedge});
\item beyond squarefree $\chi_w$: degeneracy strata and a Kummer
      dichotomy for prime $q$ (\Cref{thm:D2}, \S\ref{sec:degeneracy});
\item beyond scalar operators: block-diagonal systems and a symmetric-power
      target (\Cref{thm:S1}, \S\ref{sec:systems});
\item beyond the discrete invariant: an analytic refinement attaching
      anti-Stokes directions to differences $\Dij = c_i - c_j$
      (\Cref{thm:AT1,thm:AT3,thm:AT5,thm:AT7}, \S\ref{sec:analytic}).
\end{itemize}

\subsection{What this is not}\label{ssec:intro-scope}

We do not classify operators, and we make no claim that $\Cl(L)$ is a
complete formal invariant: it is a coarsening of the irregular-type data
of Levelt--Turrittin, and the kernel of that coarsening (operators with the
same constellation but different Stokes structures) is itself interesting
and remains open. \Cref{sec:related} discusses how the present
classification relates to four neighbouring programmes---ADE
singularities; Stokes-graph / spectral-network classifications; the
formal classification of Sabbah--Mochizuki and Birkhoff--Trjitzinsky;
wild character varieties and irregular Riemann--Hilbert---and explains in
what precise sense our object is orthogonal to each.

\subsection{Organisation}\label{ssec:intro-organisation}

\S\ref{sec:prelim} fixes notation and recalls Newton polygons. \S\ref{sec:main}
proves the Main Theorem. \S\ref{sec:multiedge} extends to multi-edge polygons
(M1) and nested ramification (N1). \S\ref{sec:degeneracy} addresses the
degeneracy strata, the Vieta monodromy lift, and the Kummer dichotomy (D2).
\S\ref{sec:systems} treats systems (S1). \S\ref{sec:analytic} states and
proves the four analytic-lifting theorems (AT1, AT3, AT5, AT7).
\S\ref{sec:classes} presents the $23$-class atlas, the taxonomy diagram, and
the constellation zoo. \S\ref{sec:related} positions the work next to its
neighbours, and \S\ref{sec:vision} closes with future directions.
\Cref{sec:reproducibility} contains a reproducibility note and points to
the companion dataset \cite{wkb_dataset_v1}.

\bigskip
The reader interested only in the main result can read
\Cref{sec:main} after \S\ref{sec:prelim}; the
analytic theorems of \S\ref{sec:analytic} are logically dependent on
\Cref{sec:main,sec:multiedge} but can be read in any internal order.


%% ============================================================================
%% §2  PRELIMINARIES
%% ============================================================================
\section{Newton polygons and the WKB symbol}\label{sec:prelim}

\subsection{Operators and conventions}\label{ssec:prelim-ops}

We work over $\mathbb{C}$. Let $L$ be a scalar linear operator
\[
   L \;=\; \sum_{k=0}^{n}\, a_k(x)\, \op^{\,k},
   \qquad a_k \in \mathbb{C}[x],\ a_n \ne 0,
\]
where $\op$ is one of:
\begin{itemize}[leftmargin=2em]
\item $\opD := \mathrm{d}/\mathrm{d}x$ \quad (the \emph{ODE case}),
\item $\opS \colon f(x) \mapsto f(x+1)$ \quad (the \emph{difference case}),
\item $\opT \colon f(x) \mapsto f(qx)$ for fixed $|q| \ne 0, 1$ \quad
      (the \emph{$q$-difference case}; we reserve the symbol $q$ for the
      slope denominator below and use a separate parameter for the
      $q$-shift when both appear).
\end{itemize}
We assume $L$ has a single irregular singularity at $x = \infty$ and is
regular elsewhere. The coefficients $a_k$ are polynomial throughout; the
extension to formal Laurent series is straightforward and notationally
heavier without affecting the structural law.

\subsection{Newton polygon at infinity}\label{ssec:prelim-Newton}

For each $a_k \in \mathbb{C}[x] \setminus \{0\}$, write $\deg a_k$ for its
degree in $x$. The \emph{Newton polygon at infinity} of $L$ is
\[
   \NL \;:=\;
   \mathrm{ConvexHull}\big\{\,(k, \deg a_k) : a_k \ne 0\,\big\}
   \;-\; \mathrm{lower\ half\!-plane},
\]
\ie\ the upper convex hull of the support, viewed as a polygon in
$\mathbb{R}^2$. Edges of $\NL$ are line segments with slope $-p/q$ in
lowest terms ($\gcd(p, q) = 1$, $p, q \ge 1$), oriented so that the slope
decreases left-to-right.

\begin{definition}[Dominant edge]\label{def:dominant-edge}
A \emph{dominant edge} of $\NL$ is an edge of strictly maximal slope. We
say $L$ has \emph{one dominant edge} if $\NL$ admits a unique such edge.
Multi-edge polygons are treated in \S\ref{sec:multiedge}.
\end{definition}

\begin{remark}[Equation-type independence of $N(L)$]\label{rem:Newton-uniform}
The Newton polygon is defined uniformly across $\opD$, $\opS$, $\opT$:
the recipe ``plot $(k, \deg a_k)$ and take the upper convex hull'' produces
the same polygon. The \emph{analytic} interpretation of an edge differs
(slope $-p/q$ means $f \sim \exp(c\, x^{p/q})$ for $\opD$, but
$f \sim \prod (\text{shift factors})$ for $\opS$ or $\opT$); the
\emph{combinatorial} content is identical.
\end{remark}

\begin{lemma}[Lattice points on an edge]\label{lem:lattice-on-edge}
Let $E$ be an edge of $\NL$ of slope $-p/q$ with $\gcd(p, q) = 1$. Then
the lattice points of $E$ are exactly
\[
   E \cap \mathbb{Z}^2
   \;=\;
   \big\{(k_0 + jq,\ h_0 - jp) : 0 \le j \le r\big\},
\]
where $(k_0, h_0)$ is the lower-right endpoint of $E$ and $r \ge 1$ is the
\emph{rank multiplier}. Equivalently, $E$ contains $r + 1$ lattice points
and has horizontal extent $rq$.
\end{lemma}

\begin{proof}
A vector $(\Delta k, \Delta h)$ along $E$ satisfies
$\Delta h / \Delta k = -p/q$. The primitive integer vector in this
direction is $(q, -p)$, and $\gcd(q, p) = 1$ ensures these are the only
lattice points lying on $E$ between its two endpoints, parametrised by
$j = 0, 1, \dots, r$ where $r$ is the integer length of $E$ measured in
multiples of $(q, -p)$.
\end{proof}

\subsection{The WKB symbol and the characteristic polynomial}\label{ssec:prelim-WKB}

Fix a dominant edge $E$ of slope $-p/q$ with lattice points
$(k_0 + jq,\ h_0 - jp)$, $0 \le j \le r$, as above. For each such lattice
point, write
\[
   \alpha_j
   \;:=\;
   \LCinf\big(a_{k_0 + jq}\big)
   \;\in\; \mathbb{C},
\]
where $\LCinf$ extracts the leading coefficient in $x$ at $\infty$.
(The lattice point lies on $E$ iff $a_{k_0 + jq}$ has degree exactly
$h_0 - jp$ in $x$, so $\alpha_j$ is unambiguous.)

\begin{definition}[WKB characteristic polynomial]\label{def:chi}
The \emph{WKB characteristic polynomial} of $L$ along the dominant edge $E$
is
\[
   \chi(c)
   \;:=\;
   \sum_{j=0}^{r}\, \alpha_j\, c^{(r-j)q}
   \;\in\; \mathbb{C}[c].
\]
The \emph{inner polynomial} is
\[
   \chi_w(w)
   \;:=\;
   \sum_{j=0}^{r}\, \alpha_j\, w^{\,r-j}
   \;\in\; \mathbb{C}[w],
   \qquad \deg \chi_w = r.
\]
The \emph{constellation} of $L$ is the multiset
$\Cl(L) := \Roots(\chi) \subset \mathbb{C}$, with roots counted with
multiplicity.
\end{definition}

\begin{remark}[Origin of $\chi$ from the WKB ansatz]\label{rem:chi-from-ansatz}
For $\op = \opD$, substituting $f = \exp\!\big(\frac{q}{p}\, c\, x^{p/q}\big)\,
g(x)$ into $L f = 0$ with $c$ to be determined gives, at leading order in
$x$, the equation $\sum_j \alpha_j\, c^{(r-j)q}\, x^{(r-j)p + h_0} = 0$ along
the edge $E$. Equating coefficients of the highest power of $x$ yields
$\chi(c) = 0$, recovering \Cref{def:chi}. For $\opS$, the substitution is
$f \sim \exp\!\big(\frac{c\, x^{1 + p/q}}{1 + p/q}\big)$ and the algebra is
identical; likewise for $\opT$ (with the appropriate logarithm). The
\emph{combinatorial} object $\chi$ is identical in all three cases; this
is the source of equation-type independence at the constellation level.
\end{remark}

\begin{notation}[The classification functor]\label{not:F}
For brevity we package $\Fclass \colon L \mapsto \Cl(L)$ as a map from the
category $\mathsf{Op}$ of scalar local irregular operators (up to formal
gauge fixing the leading edge) to the category $\mathsf{MSet}_\mathbb{C}$
of finite multisets in $\mathbb{C}$ modulo global rotation. ``Functor'' is
used here in the loose sense: the assignment respects formal gauge
equivalence; functoriality with respect to morphisms in $\mathsf{Op}$
(\eg\ pullback along a finite ramification, twist by an exponential)
is recorded as needed below.
\end{notation}

%% ============================================================================
%% §3  MAIN THEOREM
%% ============================================================================
\section{The Main Theorem}\label{sec:main}

We now state and prove the structural law that organises the rest of the
paper. Hypotheses are kept minimal: a scalar operator with one dominant
edge, in any of the three equation types.

\begin{theorem}[Factorisation of the WKB characteristic polynomial]\label{thm:main}
Let $L = \sum_{k=0}^{n} a_k(x)\, \op^{\,k}$ be a scalar linear operator with
polynomial coefficients $a_k \in \mathbb{C}[x]$ and a single irregular
singularity at $x = \infty$, where $\op \in \{\opD, \opS, \opT\}$. Assume
the Newton polygon $\NL$ at $\infty$ has a unique dominant edge $E$ of
slope $-p/q$ with $\gcd(p, q) = 1$ and $p, q \ge 1$, containing $r + 1$
lattice points
\[
   E \cap \mathbb{Z}^2 \;=\;
   \big\{(k_0 + jq,\ h_0 - jp) : 0 \le j \le r\big\}.
\]
Set $\alpha_j := \LCinf(a_{k_0 + jq})$ for $0 \le j \le r$, and define
$\chi(c)$ and $\chi_w(w)$ as in \Cref{def:chi}. Then
\[
   \boxed{\,
     \chi(c) \;=\; \chi_w(c^q),
   \,}
\]
and the WKB exponent constellation factors as the $\mu_q$-equivariant
fibre product
\[
   \Cl(L) \;=\; \Roots(\chi)
   \;=\;
   \bigsqcup_{w_i \in \Roots(\chi_w)}
   \big\{\, c \in \mathbb{C} : c^q = w_i \,\big\},
\]
where roots are counted with multiplicity and the union is over the $r$
roots $w_i$ of $\chi_w$.
\end{theorem}

\begin{proof}
By \Cref{def:chi},
$\chi(c) = \sum_{j=0}^{r} \alpha_j\, c^{(r-j) q}
       = \sum_{j=0}^{r} \alpha_j\, (c^q)^{r-j}
       = \chi_w(c^q).$
The factorisation of $\Cl(L)$ follows by substitution: $\chi(c) = 0$ iff
$c^q \in \Roots(\chi_w)$, and for each such $w_i$ the preimage under
$c \mapsto c^q$ is the $\mu_q$-orbit $\{c : c^q = w_i\}$. Multiplicities
in $\chi_w$ are transmitted to $\chi$ unchanged: if $w_i$ is a root of
$\chi_w$ of multiplicity $m_i$, then each $c$ with $c^q = w_i$ is a root
of $\chi$ of multiplicity $m_i$.
\end{proof}

\begin{corollary}[Generic ring decomposition]\label{cor:ring}
If the moduli $|w_1|, \dots, |w_r|$ are pairwise distinct and all
$w_i \ne 0$, then $\Cl(L)$ consists of exactly $r$ concentric
$\mu_q$-orbits (``rings'') of $q$ points each, with radii $|w_i|^{1/q}$
and angular bases $\arginline(w_i) / q \pmod{2\pi/q}$.
\end{corollary}

\begin{proof}
For each simple root $w_i$ of $\chi_w$, the fibre $\{c : c^q = w_i\}$
consists of $q$ points evenly spaced on the circle of radius
$|w_i|^{1/q}$, with one of them at angle $\arginline(w_i)/q$ and the others
obtained by multiplication by $\zeta_q^k$. The hypotheses make the
$r$ rings concentric and pairwise distinct.
\end{proof}

\begin{corollary}[Vieta monodromy]\label{cor:vieta}
Fix $\alpha_0$ and let $\beta := \alpha_r$ vary. Under $\beta \mapsto
e^{i\theta} \beta$,
\[
   \arginline\!\Big(\prod_{c \in \Cl(L)} c\Big)
   \;\longmapsto\;
   \arginline\!\Big(\prod_{c \in \Cl(L)} c\Big) \;+\; \theta \pmod{2\pi}.
\]
If $\Cl(L)$ admits a $\mu_q$-equivariant ring decomposition
$\Cl(L) = \bigsqcup_{k=1}^{s} R_k$ with $|R_k| = m_k$, and the rings deform
rigidly under the variation of $\beta$ with rotations $\Delta_k$ (each
defined modulo $2\pi / m_k$), then the unambiguous lifts $m_k \Delta_k \in
\mathbb{R}/2\pi\mathbb{Z}$ obey
\[
   \sum_{k=1}^{s} m_k\, \Delta_k \;\equiv\; \theta \pmod{2\pi}.
\]
\end{corollary}

\begin{proof}
Vieta's formula gives $\prod_{c \in \Roots(\chi)} c = (-1)^{\deg \chi}\,
\alpha_0 / \alpha_r$, so its argument is shifted by $-\theta$ under
$\alpha_r \mapsto e^{i\theta} \alpha_r$ (sign convention absorbed). For the
ring lift: each ring $R_k$ contributes $m_k \Delta_k$ to the change in
$\arginline \prod c$ (the per-ring rotations are well-defined modulo
$2\pi/m_k$, and the lift $m_k \Delta_k$ is unambiguous). Summing and
matching to the total shift $\theta$ gives the rule.
\end{proof}

\begin{remark}[Degenerate loci]\label{rem:degenerate}
The clean ring picture of \Cref{cor:ring} fails on three explicit
subvarieties of coefficient space $(\alpha_0, \dots, \alpha_r)$:
\begin{enumerate}[label=(\roman*)]
\item the \emph{discriminant locus} $\{\disc \chi_w = 0\}$, where roots of
      $\chi_w$ coalesce and ring multiplicities exceed $q$;
\item the \emph{equal-modulus locus} $\{|w_i| = |w_j|$ for some $i \ne j\}$,
      where distinct $\mu_q$-orbits merge into a single radius with
      non-uniform angular distribution;
\item the \emph{zero-root locus} $\{\alpha_r = 0\}$, where $c = 0$ enters
      the constellation with multiplicity $q$ and the dominant edge itself
      jumps.
\end{enumerate}
The first carries resonance / logarithmic Stokes data; the second is
responsible for the apparent regularity of stress test
\texttt{STRESS\_equalmod} (a six-point constellation on one circle that is
\emph{not} a regular hexagon); the third is a wall-crossing of the Newton
polygon itself. \S\ref{sec:degeneracy} treats them systematically.
\end{remark}

\begin{remark}[Equation-type invariance]\label{rem:eq-invariance}
The proof of \Cref{thm:main} uses only the combinatorial identity
$c^{(r-j)q} = (c^q)^{r-j}$, no analytic input. Consequently, the
factorisation holds identically across $\opD$, $\opS$, $\opT$. The analytic
\emph{meaning} of $c \in \Cl(L)$ differs (it is the leading-order WKB
exponent $c$ in $f \sim \exp(c\, x^{p/q}/(p/q))$ for $\opD$; for $\opS$ it
is $\log(S f / f)|_\infty$ to leading order; for $\opT$ it is
$\log_q(T f / f)|_\infty$); the multiset $\Cl(L)$ is invariant. This is
the precise sense in which $\Fclass$ is equation-type independent.
\end{remark}

\begin{remark}[Forward references]\label{rem:forward-main}
\Cref{thm:main} is the base case of the entire programme. It generalises
in three directions: to multi-edge polygons (\Cref{thm:M1}); to nested
factorisations of $\chi_w$ (\Cref{thm:N1}); to systems (\Cref{thm:S1}).
The complementary analytic refinement
(\Cref{thm:AT1,thm:AT3,thm:AT5,thm:AT7}) attaches anti-Stokes directions
to the differences $\Dij = c_i - c_j$ on the ramified cover $x = t^q$.
\end{remark}


%% ============================================================================
%% §4  MULTI-EDGE AND NESTED RAMIFICATION
%% ============================================================================
\section{Multi-edge polygons and nested ramification}\label{sec:multiedge}

\subsection{Per-edge factorisation along the slope filtration}\label{ssec:M1}

When $N(L)$ has several edges of distinct slopes, the Main Theorem applies
per edge and the constellations assemble multiplicatively via the
formal-classification (Levelt--Turrittin) slope decomposition.

\begin{theorem}[Per-edge factorisation; multi-edge generalisation]\label{thm:M1}
Let $L$ be a scalar linear operator with polynomial coefficients on
$\mathbb{C}[x]$ and a single irregular singularity at $x = \infty$. Suppose
$\NL$ has edges $E_1, \dots, E_m$ with pairwise distinct slopes
$-p_\alpha / q_\alpha$ ($\gcd(p_\alpha, q_\alpha) = 1$) and $r_\alpha + 1$
lattice points on $E_\alpha$. Let $k_0^{(\alpha)}$ be the height of $E_\alpha$
above the $c$-axis, and let $\chi_w^{(\alpha)} \in \mathbb{C}[w]$ be the
inner polynomial of degree $r_\alpha$ built from the $\LCinf$ data of
$E_\alpha$ (per \Cref{def:chi}). Then:
\begin{enumerate}[leftmargin=2em]
\item (Slope decomposition.) The formal differential module
      $\widehat{M}(L)$ associated to $L$ at $x = \infty$ admits a canonical
      direct-sum decomposition $\widehat{M}(L) \simeq \bigoplus_{\alpha=1}^{m}
      \widehat{M}_\alpha$ with $\widehat{M}_\alpha$ of pure slope
      $p_\alpha / q_\alpha$. Choosing cyclic vectors realises each
      $\widehat{M}_\alpha$ as the formal module of a scalar single-slope
      operator $L_\alpha$, and we write
      $\widehat{L} \simeq L_1 \oplus \cdots \oplus L_m$.
\item (Per-edge characteristic polynomial.) Each $L_\alpha$ satisfies
      \Cref{thm:main} with
      $\chi^{(\alpha)}(c) = c^{k_0^{(\alpha)}}\, \chi_w^{(\alpha)}(c^{q_\alpha})$.
\item (Total factorisation.) Setting $\chi(c) := \prod_{\alpha=1}^{m}
      \chi^{(\alpha)}(c)$, one has
      $\Cl(L) = \bigsqcup_{\alpha=1}^{m} \Cl_\alpha$ as multisets, with
      $\Cl_\alpha := \Roots(\chi^{(\alpha)})$.
\end{enumerate}
On the open subset $U_{\mathrm{disj}}$ of coefficient space where the
per-edge radial annuli $[\rho_\alpha^{\min}, \rho_\alpha^{\max}]$ with
$\rho_\alpha^{\bullet} = |w^\bullet|^{1/q_\alpha}$ are pairwise disjoint,
the union is a disjoint union of \emph{distinct} subsets of $\mathbb{C}$.
\end{theorem}

\begin{proof}[Proof sketch]
\textit{Step 1 (Slope decomposition).} By the Levelt--Turrittin theorem
(Sabbah \cite[Ch.~II]{sabbah-isomon}; Wasow \cite[Ch.~VII]{wasow};
Sibuya \cite{sibuya}) the formal differential module $\widehat{M}(L)$
decomposes canonically as a direct sum of slope-graded pieces, with the
slopes of the summands being exactly the slopes of the edges of $\NL$.
Choosing cyclic vectors in each summand realises $\widehat{M}_\alpha$ as
the formal module of a single-slope scalar operator $L_\alpha$, well
defined up to formal gauge.

\textit{Step 2 (Per-edge Main Theorem).} Each $L_\alpha$ is a scalar
single-slope connection with Newton polygon a single edge of slope
$-p_\alpha / q_\alpha$ and $r_\alpha + 1$ lattice points; \Cref{thm:main}
applies and gives
$\chi^{(\alpha)}(c) = c^{k_0^{(\alpha)}}\, \chi_w^{(\alpha)}(c^{q_\alpha})$,
where $k_0^{(\alpha)}$ records the multiplicity of $c = 0$ contributed by
lattice points strictly below $E_\alpha$.

\textit{Step 3 (Product factorisation).} The WKB characteristic polynomial
of a direct sum of formal connections is the product of the per-summand
characteristic polynomials (multiplicativity of the determinant of the
horizontal connection $1$-form, equivalently of the Newton polygon under
$\oplus$; \cf\ Robba--Malgrange \cite{robba-malgrange}). Hence
$\chi(c) = \prod_\alpha \chi^{(\alpha)}(c)$ and $\Cl(L) = \bigsqcup_\alpha
\Cl_\alpha$ as multisets.

\textit{Step 4 (Genericity).} On the open subset $U_{\mathrm{disj}}$ the
per-edge constellations occupy disjoint radial annuli, so the multiset
union is an honest disjoint union of subsets of $\mathbb{C}^\times$. On
the complement (radii-collision stratum), the unions remain disjoint as
multisets, with possible coincidences in $\mathbb{C}^\times$ recorded by
multiplicity. \qedhere
\end{proof}

\begin{remark}[Counterexample boundary for M1]\label{rem:M1-counter}
The disjoint-union conclusion in (3) is a multiset statement; as
\emph{sets} in $\mathbb{C}^\times$ the per-edge constellations may
coincide on the \emph{inter-edge equal-radius stratum} (examples in the
companion dataset: \texttt{MULTI\_edges\_1o4\_then\_1o2},
\texttt{MULTI\_edges\_1o3\_r2\_then\_1o2}; see
\cite[counterexamples.json: CE-multiedge-equalradius]{wkb_dataset_v1}).
Likewise, when a per-edge $\chi_w^{(\alpha)}(0) = 0$, the zero-root
contributions of adjacent edges \emph{share} the origin in $\mathbb{C}$
(CE-multiedge-zerosharing). These are not contradictions: the multiset
disjoint union is preserved; only the set-theoretic statement is
weakened.
\end{remark}

\subsection{Nested ramification}\label{ssec:N1}

Beyond a single $\mu_q$-pullback, the inner polynomial $\chi_w$ itself
may factor through a higher power.

\begin{theorem}[Nested ramification: iterated $\mu_q$-pullback]\label{thm:N1}
Let $L$ be scalar with one dominant edge of slope $-p/q$ as in
\Cref{thm:main}. Suppose the inner polynomial admits a factorisation
$\chi_w(w) = \chi_v(w^{q'})$ for some integer $q' \ge 2$ and a polynomial
$\chi_v \in \mathbb{C}[v]$ of degree $r' = r / q'$ (\ie\ $\chi_w$ has
nonzero coefficients only in degrees that are multiples of $q'$). Then
\[
   \Cl(L) \;=\; \bigsqcup_{v \in \Roots(\chi_v)}
            \big\{\, c \in \mathbb{C} : c^{qq'} = v \,\big\}.
\]
Each $v$-fibre $\{c : c^{qq'} = v\}$ has size $qq'$ and carries a free
$\mu_{qq'}$-action. The \emph{visible} symmetry group generated by the
outer $\mu_q$-action (from $\chi(c) = \chi_w(c^q)$) and the nested
$\mu_{q'}$-action (from $\chi_w = \chi_v \circ (\cdot)^{q'}$) is
\[
   \mu_{\mathrm{lcm}(q, q')} \;\subseteq\; \mu_{qq'},
\]
with equality $\mu_{\mathrm{lcm}(q, q')} = \mu_{qq'}$ if and only if
$\gcd(q, q') = 1$ (in which case $\mu_{qq'} \cong \mu_q \times \mu_{q'}$
canonically by CRT).
\end{theorem}

\begin{proof}
By \Cref{thm:main}, $\chi(c) = \chi_w(c^q)$; substituting the hypothesis
gives $\chi(c) = \chi_v\big((c^q)^{q'}\big) = \chi_v(c^{qq'})$. The fibre
decomposition follows. For the symmetry claim, the original Main
Theorem's $\mu_q$-action sits inside $\mu_{qq'}$ as the subgroup
$\{\zeta_{qq'}^{q' k}\}$. A lift to $c$-level of the inner-polynomial
$\mu_{q'}$-symmetry (acting on $w = c^q$ by $w \mapsto \zeta_{q'} w$) is
multiplication by an $\eta \in \mu_{qq'}$ with $\eta^q = \zeta_{q'}$;
such $\eta$ exists (\eg\ $\eta = \zeta_{qq'}^q$) and is unique modulo
$\mu_q$. The subgroup of $\mu_{qq'}$ generated by $\zeta_q$ and $\eta$
is $\mu_{\mathrm{lcm}(q, q')}$.
\end{proof}

\begin{remark}[Minimality]\label{rem:N1-min}
The hypothesis $\chi_w(w) = \chi_v(w^{q'})$ is preserved under further
inner factorisations $\chi_v(v) = \chi_u(v^{q''})$; iterating until no
inner factorisation is possible gives the \emph{minimal} $q' q'' \cdots$
through which the constellation factors. The Phase-III sweep records two
such nested operators in the atlas (\texttt{NEST\_q3\_qinner2},
\texttt{NEST\_q2\_qinner3}), both stable under one round of inner
factorisation.
\end{remark}

\begin{remark}[Constellation vs.\ analytic ramification]\label{rem:N1-analytic}
\Cref{thm:N1} is a purely algebraic identity at the constellation level.
The analytic implications for the Levelt--Turrittin ramification cover are
deferred to \Cref{thm:AT3}; crucially, the minimal cover for the formal
classification remains $x = t^q$ and does not increase to $x = \tau^{qq'}$.
The nested factorisation refines the wild character variety, not the
ramification index.
\end{remark}


%% ============================================================================
%% §5  DEGENERACY STRATA AND ARITHMETIC STRATIFICATION
%% ============================================================================
\section{Degeneracy strata and arithmetic stratification}\label{sec:degeneracy}

The Main Theorem (\Cref{thm:main}) decomposes the image of $\Fclass$
into a generic stratum and the three degenerate loci of
\Cref{rem:degenerate}. We now treat the degeneracy strata and their
arithmetic refinement, culminating in a prime-$q$ Kummer dichotomy
(\Cref{thm:D2}).

\subsection{The four named strata}\label{ssec:strata-named}

In the companion atlas \cite{wkb_dataset_v1}, four stress-test operators
exhibit the qualitatively distinct degenerate behaviours:
\begin{enumerate}[leftmargin=2em]
\item \texttt{STRESS\_equalmod} (equal-modulus locus): six points on a
      single circle of radius $\sqrt{2}$, \emph{not} evenly spaced.
\item \texttt{STRESS\_doubleroot} (discriminant locus / resonant
      multiplicity): two rings, the inner with multiplicity $2$ at each
      apparent point.
\item \texttt{STRESS\_zeroroot} (zero-root locus / wall-crossing): two
      points at radius $1$ together with a multiplicity-$q$ root at the
      origin.
\item \texttt{STRESS\_accidental\_sym} (enhanced symmetry): a regular
      $rq$-gon with cyclic symmetry $C_{rq} \supsetneq C_q$, occurring at
      slope $-p/q$ but exhibiting symmetry above what $q$ predicts.
\end{enumerate}
The Vieta lift (\Cref{cor:vieta}) characterises the equal-modulus stratum
behaviour (rigid permutation up to shared rotation, not rigid rotation).
The discriminant locus carries logarithmic / resonant Stokes data
(\Cref{rem:degenerate}(i)). The zero-root locus is a wall-crossing of
$\NL$ itself: when $\alpha_r = 0$ the dominant edge has been mis-identified
and is replaced by a strictly shorter edge with smaller $r$.

\subsection{\texorpdfstring{The Kummer dichotomy for prime $q$}{The Kummer dichotomy for prime q}}\label{ssec:D2}

Beyond the open stratum, the arithmetic of the inner polynomial controls
whether the per-ring constellations admit nontrivial integer linear
relations. The classical setting is Kummer theory: when ratios of inner
roots are $q$-th powers in $\mathbb{Q}(\zeta_q)$, the constellation
collapses arithmetically.

\begin{theorem}[Inner-polynomial $\mathbb{Z}$-resonance dichotomy; prime
$q$]\label{thm:D2}
Let $L$ be scalar with one dominant edge of slope $-p/q$, $q$ \emph{prime},
and let $\chi_w$ be squarefree with $r \ge 2$ distinct roots
$w_1, \dots, w_r$. Choose $q$-th roots $z_\alpha = w_\alpha^{1/q}
\in \mathbb{C}$ so that
$\Cl(L) = \bigsqcup_\alpha R_\alpha$ with
$R_\alpha = \{\zeta_q^j z_\alpha : 0 \le j < q\}$.
For $\lambda_1, \dots, \lambda_r \in \Zlc$, set $\Lambda(\lambda) =
\sum_\alpha \lambda_\alpha\, z_\alpha \in \mathbb{C}$. Call a
$\mathbb{Z}$-relation $\sum_{\alpha, j} n_{\alpha j}\, \zeta_q^j\, z_\alpha
= 0$ \emph{$\mu_q$-forced} if every $\lambda_\alpha := \sum_j n_{\alpha j}
\zeta_q^j$ vanishes in $\Zlc$; \emph{cross-ring nontrivial} otherwise.
Then:
\begin{enumerate}[leftmargin=2em]
\item \textbf{Generic stratum.} If $\{w_1, \dots, w_r\}$ are algebraically
      independent over $\mathbb{Q}$, then $\Cl(L)$ admits \emph{no}
      cross-ring nontrivial $\mathbb{Z}$-relation.
\item \textbf{Kummer/cyclotomic stratum.} If there exist
      $\lambda_\alpha \in \Zlc$ not all zero with $\sum_\alpha
      \lambda_\alpha\, z_\alpha = 0$ in $\mathbb{C}$, then this identity
      \emph{is} a cross-ring nontrivial $\mathbb{Z}$-relation. A clean
      sufficient condition is $w_\beta / w_\alpha \in (\Qlc^\times)^q$ for
      some $\alpha \ne \beta$ (so $z_\beta / z_\alpha \in \Qlc$ after
      consistent choices of $q$-th roots).
\end{enumerate}
\end{theorem}

\begin{proof}[Proof sketch]
\textit{Set-up.} Each $R_\alpha = \{\zeta_q^j\, z_\alpha\}_{j=0}^{q-1}$.
Re-grouping by $\alpha$, the relation $\sum_{\alpha, j} n_{\alpha j}\,
\zeta_q^j\, z_\alpha = 0$ becomes $\sum_\alpha \lambda_\alpha\, z_\alpha
= 0$ in $\mathbb{C}$, with $\lambda_\alpha \in \Zlc$.

\textit{$\mu_q$-forced sub-lattice.} For prime $q$, $\Zlc \cong \mathbb{Z}[x]
/ \Phi_q(x)$ has rank $q - 1$, so the kernel of $\mathbb{Z}^q \to \Zlc$,
$(n_0, \dots, n_{q-1}) \mapsto \sum_j n_j\, \zeta_q^j$, is generated by
the cyclotomic relation $1 + \zeta_q + \cdots + \zeta_q^{q-1} = 0$. Hence
$\lambda_\alpha = 0$ in $\Zlc$ iff $(n_{\alpha 0}, \dots, n_{\alpha, q-1})$
is a constant vector; in particular the only $\mu_q$-forced relation on a
single ring is the per-ring sum $\sum_{c \in R_\alpha} c = 0$.

\textit{Case (i): Algebraic independence.} If $\{w_\alpha\}$ are
algebraically independent over $\mathbb{Q}$, then $\{z_\alpha\}$ are
algebraically independent over $\Qlc$ (Kummer/Eisenstein: each $z_\alpha$
is algebraic of degree $q$ over $\Qlc(w_\alpha)$, and the $w_\alpha$ are
independent). Hence any $\Zlc$-linear dependence $\sum_\alpha \lambda_\alpha
z_\alpha = 0$ forces all $\lambda_\alpha = 0$, \ie\ the original
$\mathbb{Z}$-relation is $\mu_q$-forced.

\textit{Case (ii): Kummer / cyclotomic.} If some $\lambda_\alpha \ne 0$
exists with $\sum \lambda_\alpha z_\alpha = 0$, expanding $\lambda_\alpha
= \sum_j n_{\alpha j}\, \zeta_q^j$ recovers a $\mathbb{Z}$-relation not in
the $\mu_q$-forced sub-lattice; it is cross-ring nontrivial. The
sufficient conditions follow from Kummer theory: if $w_\beta / w_\alpha
= u^q$ with $u \in \Qlc^\times$, choose $z_\beta = u\, z_\alpha$, giving
the relation $z_\beta - u\, z_\alpha = 0$. Specialising to $u \in
\mathbb{Q}^\times$ gives a $\mathbb{Z}$-relation with support entirely on
$\{z_\alpha, z_\beta\}$.

\textit{Sharpness on the atlas.} Every prime-$q$ generic-$\chi_w$ operator
in the Phase-III sweep ($q \in \{5, 7\}$) produces zero cross-ring
$\mathbb{Z}$-relations in the search window $|n_k| \le 3$, support
$\le 4$; every simple-$\chi_w$ multi-ring operator produces $\ge 5$. The
dichotomy is exact for prime $q$ within the search window.
\end{proof}

\begin{remark}[Composite $q$]\label{rem:D2-composite}
For composite $q$, $\Zlc$ has further structure: the kernel of the
projection $\Zlc \to \prod_{d \mid q} \mathbb{Z}[\zeta_d]$ captures
sub-cyclotomic resonances, and the $\mu_q$-forced sub-lattice includes
sub-orbit sums coming from each divisor $d \mid q$. The corresponding
dichotomy holds with ``$\lambda_\alpha = 0$ in $\Zlc$'' replaced by
``$\lambda_\alpha$ in the kernel of the projection''. Examples implicit
in $q = 4$ and $q = 6$ operators of the atlas illustrate the refinement.
\end{remark}

\begin{remark}[Search-window limitation]\label{rem:D2-search}
\Cref{thm:D2}(i) is stated within the Phase-III search window
($|n_k| \le 3$, support $\le 4$); the genuine algebraic-independence
statement (no $\mathbb{Z}$-relation \emph{at all}) is a
transcendental-degree question whose verification requires reasoning
beyond a bounded search. The companion counterexample knowledge base
\cite[CE-Zres-largecoef]{wkb_dataset_v1} records the boundary.
\end{remark}


%% ============================================================================
%% §6  SYSTEMS
%% ============================================================================
\section{Systems}\label{sec:systems}

The constellation functor extends to systems whose leading symbol is
semisimple, via a block-diagonal pullback.

\begin{theorem}[System pullback: block-diagonal direct sum]\label{thm:S1}
Let $\Phi = \mathrm{diag}(L_1, \dots, L_n)$ be a block-diagonal system of
scalar linear operators, each $L_i$ scalar with a single irregular
singularity at $x = \infty$ and Newton polygon $N(L_i)$. Then
\[
   \Cl(\Phi) \;=\; \bigsqcup_{i = 1}^{n}\, \Cl(L_i)
   \qquad \text{as multisets,}
\]
and the constellation functor factors
\[
   \Fclass \colon \Sys^n(\mathrm{scalar}) \;\longrightarrow\;
   \Sym^n(\text{finite multisets in } \mathbb{C}),
   \qquad \Fclass(\Phi) \;=\; \bigsqcup_i \Fclass(L_i).
\]
\end{theorem}

\begin{proof}
For a block-diagonal system, the WKB symbol of $\Phi$ acts diagonally on
each block, so the leading balance decouples and the WKB characteristic
polynomial of $\Phi$ is $\chi_\Phi(c) = \prod_i \chi_{L_i}(c)$. Hence
$\Roots(\chi_\Phi) = \bigsqcup_i \Roots(\chi_{L_i})$ as multisets. The
factorisation through the symmetric power is the observation that the
constellation is a multiset, not an ordered tuple, so the assignment
$\Phi \mapsto (\Cl(L_1), \dots, \Cl(L_n))$ is $\Symm_n$-equivariant.
\end{proof}

\begin{remark}[Diagonalisable coupled systems]\label{rem:S1-diag}
For a general first-order system whose leading symbol $A_0(x) =
\LCinf(\Phi)$ is diagonalisable with distinct eigenvalues, a formal gauge
reduces $\Phi$ to block-diagonal form to leading order, and
\Cref{thm:S1} applies at the constellation level. Coalescence of leading
eigenvalues defines a \emph{system discriminant stratum} where the
Levelt--Turrittin classification of irregular formal types is required;
the present theorem does not extend to that stratum without further
hypotheses, and genuinely coupled systems (non-semisimple leading symbol,
Jordan blocks) lie outside the present scope. See Boalch
\cite{boalch-stokes,boalch-wildhodge} for the moduli-theoretic treatment.
\end{remark}

\begin{remark}[Counterexample boundary for S1]\label{rem:S1-counter}
The companion counterexample catalogue (CE-coupled-system) records that
generic Painlev{\'e}-type and coupled-system examples lie outside the
block-diagonal hypothesis; in such cases the leading-symbol diagonalisation
is itself a non-trivial gauge problem whose obstruction sits in the
system-discriminant stratum.
\end{remark}


%% ============================================================================
%% §7  ANALYTIC LIFTING
%% ============================================================================
\section{Analytic lifting: anti-Stokes skeleton and slope filtration}\label{sec:analytic}

\subsection{Open-stratum hypotheses}\label{ssec:OS}

The analytic theorems below organise the Stokes geometry attached to a
constellation. They are clearest on the \emph{open stratum} where no
degeneracies of \Cref{rem:degenerate} occur. We collect the running
hypotheses:
\begin{itemize}[leftmargin=2.5em]
\item[(OS1)] $\chi_w$ squarefree ($\disc \chi_w \ne 0$).
\item[(OS2)] $\chi_w(0) \ne 0$ (no zero roots in $\Cl(L)$).
\item[(OS3)] No $\omega$-resonance beyond $\mu_q$: the multiset
       $\{\arginline \Dij \mod \pi/p\}$ has no coincidences beyond
       those forced by the ring structure.
\item[(OS4)] (Multi-edge.) Per-edge radial annuli pairwise disjoint.
\item[(OS5)] (Multi-edge.) Cross-edge leading coefficients nonzero
       ($\chi_w^{(\beta)}(0) \ne 0$ for each $\beta$).
\end{itemize}
\Cref{thm:AT1} uses (OS1)--(OS3); \Cref{thm:AT3} uses (OS1)--(OS3);
\Cref{thm:AT5} uses (OS1)--(OS5); \Cref{thm:AT7} uses (OS1)--(OS3).

\subsection{\texorpdfstring{Anti-Stokes directions and $\arg \Delta_{ij}$}{Anti-Stokes directions and arg Delta\_ij}}\label{ssec:AT1}

\begin{theorem}[Anti-Stokes $\leftrightarrow$ $\arginline \Dij$ on the
cover; open stratum]\label{thm:AT1}
Let $L$ be scalar with one dominant edge of slope $-p/q$, $\op = \opD$, in
the open stratum (OS1)--(OS3). Let $\Cl(L) = \{c_1, \dots, c_N\}$ with
$N = qr$ and $\Dij = c_i - c_j$. On the ramified cover $x = t^q$, the
formal connection splits as a direct sum
\[
   \widehat{L}_{|x = t^q}
   \;\simeq\;
   \bigoplus_{i = 1}^{N}\,
   \big(\mathbb{C}\{t\},\ \mathrm{d} + \mathrm{d} Q_i(t)\big),
   \qquad Q_i(t) = c_i\, t^p + R_i(t),\ \deg R_i < p,
\]
and the anti-Stokes directions on the $t$-plane are
\[
   \Acov
   \;=\;
   \big\{\, \phi \in [0, 2\pi)\, :\,
   \arginline(\Dij) + p\, \phi \equiv \tfrac{\pi}{2} \pmod{\pi},\
   \text{ some unordered } \{i, j\} \,\big\}.
\]
Each unordered pair $\{i, j\}$ contributes exactly $2p$ cover-rays; the
total cover-ray count with multiplicity is $2 p \binom{N}{2}$. Under the
projection $t \mapsto t^q$, each pair's $2p$ cover-rays project to
$2p / \gcd(2p, q)$ base-rays.
\end{theorem}

\begin{proof}[Proof sketch]
\textit{Step 1 (Levelt--Turrittin on the $q$-cover).} By Levelt--Turrittin
(Sabbah \cite[Ch.~II]{sabbah-isomon}; Wasow \cite[Ch.~VII]{wasow};
Sibuya \cite{sibuya}), there exists a ramification $x = t^q$ on which
$\widehat{L}$ is gauge-equivalent to a direct sum of rank-one summands
with irregular type $Q_i(t) = c_i\, t^p + R_i(t)$, the leading coefficient
$c_i$ being precisely a root of $\chi$.

\textit{Step 2 (Leading anti-Stokes condition).} Anti-Stokes directions
for the pair $(i, j)$ are determined by the leading term of $Q_i - Q_j$
in $t$. On the open stratum the $c_i$ are pairwise distinct, so this
leading term is $\Dij\, t^p$. With $t = \rho\, e^{i\phi}$,
\[
   \mathrm{Re}(Q_i - Q_j)
   \;=\;
   \rho^p\, |\Dij|\, \cos(\arginline \Dij + p \phi) + O(\rho^{p-1}),
\]
and vanishing of the leading term gives $\arginline \Dij + p \phi \equiv
\pi/2 \pmod{\pi}$, with $2p$ solutions in $[0, 2\pi)$.

\textit{Step 3 (Total count).} Summing over $\binom{N}{2}$ unordered pairs
gives $2 p \binom{N}{2}$ cover-rays with multiplicity. The distinct count
is governed by collisions in $\{\arginline \Dij \mod \pi/p\}$
(\cite[phase3\_invariants.json: D1 partition]{wkb_dataset_v1}).

\textit{Step 4 (Projection to base).} The cover-to-base map $t \mapsto t^q$
acts on phases by $\phi \mapsto q\phi$; for a fixed pair the $2p$ cover
rays project to orbits of size $\gcd(2p, q)$, giving at most $2p /
\gcd(2p, q)$ base-rays per pair, with equality away from cross-pair
collisions.
\end{proof}

\begin{remark}[Phase-IV verification]\label{rem:AT1-verif}
The Phase-IV consistency check
\cite[phase4\_stokes.json]{wkb_dataset_v1} verifies the formula
$\#\Acov_{\mathrm{actual}} / \#\Acov_{\mathrm{expected}} = 1.000$ exactly
on $11$ open-stratum operators.
\end{remark}

\begin{remark}[Beyond the open stratum]\label{rem:AT1-degenerate}
Outside (OS1)--(OS3), four classes of corrections occur:
on the discriminant locus (\texttt{CTRL\_disc\_eps*}) resonance /
logarithmic Stokes data appear; on the equal-modulus locus
(\texttt{STRESS\_equalmod}, \texttt{CTRL\_conjroots\_t*}) the rotation
in \Cref{cor:vieta} becomes a rigid permutation up to shared rotation;
on the zero-root locus (\texttt{CTRL\_zeroroot\_*}) the Newton polygon
itself wall-crosses; on the $\omega$-resonance locus
(\texttt{OMEGA\_*}, \texttt{GAL\_cyclotomic\_*}) anti-Stokes rays cluster
beyond what (OS3) permits.
\end{remark}

\subsection{Nested ramification on wild data}\label{ssec:AT3}

\begin{theorem}[Nested ramification: $\mu_{q'}$-action on wild data]\label{thm:AT3}
Suppose the hypotheses of \Cref{thm:N1} hold: $\chi_w(w) = \chi_v(w^{q'})$
for some $q' \ge 2$, with $\chi_v$ admitting no further inner factorisation.
Let $\MW(L)$ denote the wild character variety of $L$ at $x = \infty$.
Then:
\begin{enumerate}[leftmargin=2em]
\item The minimal Levelt--Turrittin ramification cover for $\widehat{L}$
      is $x = t^q$ (not $x = \tau^{qq'}$): on the $q$-cover, the formal
      connection splits into rank-one summands with irregular types
      $Q_i(t) = c_i\, t^p + (\text{lower})$, and the $c_i$ are exactly
      the elements of $\Cl(L)$ listed by \Cref{thm:N1}.
\item The nested factorisation induces an \emph{additional}
      $\mu_{q'}$-action on the constellation: any $\eta \in \mu_{qq'}$
      with $\eta^q = \zeta_{q'}$ acts by $c \mapsto \eta\, c$, preserving
      each $v$-fibre $\{c : c^{qq'} = v\}$ and permuting the $q'$ choices
      of $w = c^q$ within that fibre. Combined with the cover-action
      $c \mapsto \zeta_q c$, the visible symmetry is $\mu_{\mathrm{lcm}(q,
      q')} \subseteq \mu_{qq'}$, with equality iff $\gcd(q, q') = 1$.
\item The locus of irregular types $\{Q_i(t)\}_i$ compatible with the
      factorisation $\chi_w = \chi_v \circ (\cdot)^{q'}$ is
      $\mu_{q'}$-invariant in the discrete part $\MW(L)^{\mathrm{disc}}$;
      it organises into a $\mu_{q'}$-equivariant stratum.
\end{enumerate}
\end{theorem}

\begin{proof}[Proof sketch]
\textit{Part (1).} The Levelt--Turrittin ramification index of a single-slope
scalar formal connection is the denominator $q$ of its slope in lowest
terms. The hypothesis $\chi_w(w) = \chi_v(w^{q'})$ constrains the
multiset of leading coefficients $\{c_i\}$ but does not change the slope
$p/q$ of $\NL$. The minimal cover therefore remains $x = t^q$, as in
\Cref{thm:main,thm:AT1}.

\textit{Part (2).} By \Cref{thm:N1}, $\Cl(L) = \bigsqcup_v \{c : c^{qq'}
= v\}$; each $v$-fibre carries a free $\mu_{qq'}$-action. The
$\mu_q$-action sits inside $\mu_{qq'}$ as $\{\zeta_{qq'}^{q'k}\}$, and the
$\mu_{q'}$-action of the nested factorisation lifts to a $c$-level action
by $\eta \in \mu_{qq'}$ with $\eta^q = \zeta_{q'}$. The subgroup generated
by both is $\mu_{\mathrm{lcm}(q, q')}$.

\textit{Part (3).} The discrete part $\MW(L)^{\mathrm{disc}}$
parametrises tuples $(Q_1, \dots, Q_N)$ of irregular-type polynomials
modulo formal gauge and formal monodromy. The $\mu_{q'}$-action of
Part~(2) induces a permutation $Q_i \mapsto Q_{\sigma(i)}$; the locus
compatible with a fixed $\chi_v$ is therefore $\mu_{q'}$-invariant. We
assert equivariance only; the further refinement to a $\mu_{q'}$-quotient
of an un-symmetrised moduli space lies in the wild-character-variety
formalism (Boalch \cite{boalch-stokes,boalch-wildhodge}; Sabbah
\cite{sabbah-polarizable}).
\end{proof}

\begin{remark}[Correction of a Phase-IV draft statement]\label{rem:AT3-correction}
An earlier Phase-IV draft claimed that the minimal Levelt--Turrittin cover
for a nested-ramification operator becomes $x = \tau^{qq'}$, which would
double-count anti-Stokes rays per pair. The corrected statement
(\Cref{thm:AT3}(1)) holds the cover at $x = t^q$: nested factorisation is
a refinement of the constellation and of the wild character variety
stratification, not of the formal ramification index.
\end{remark}

\subsection{Multi-edge slope filtration on the anti-Stokes skeleton}\label{ssec:AT5}

\begin{theorem}[Multi-edge slope filtration of Stokes data]\label{thm:AT5}
Suppose $L$ has multi-edge Newton polygon as in \Cref{thm:M1}, with edges
$E_\alpha$ of slope $-p_\alpha / q_\alpha$ and per-edge constellations
$\Cl_\alpha$ of size $N_\alpha = q_\alpha r_\alpha$, satisfying (OS1)--(OS5).
Let $d = \mathrm{lcm}(q_1, \dots, q_m)$ and consider the common cover
$x = t^d$, with $P_\alpha := p_\alpha d / q_\alpha$. Then:
\begin{enumerate}[leftmargin=2em]
\item (Within-edge rays.) Pairs $(i, j)$ both in $\Cl_\alpha$ contribute
      $2 P_\alpha \binom{N_\alpha}{2}$ rays (with multiplicity) at
      $t$-angle spacing $\pi / P_\alpha$.
\item (Cross-edge rays.) For pairs $(i \in \Cl_\alpha, j \in \Cl_\beta)$
      with $\alpha \ne \beta$, $P_\beta > P_\alpha$ (WLOG), the leading
      term of $Q_j^{(\beta)} - Q_i^{(\alpha)}$ on $x = t^d$ is $c_j\,
      t^{P_\beta}$ (using (OS5): $c_j \ne 0$), giving the anti-Stokes
      condition $\arginline c_j + P_\beta\, \phi \equiv \pi/2 \pmod{\pi}$
      and contributing $2 P_\beta$ rays per cross-edge pair.
\item (Slope filtration.) The Stokes data of $L$ decompose along the
      slope filtration $\widehat{L} \simeq \bigoplus_\alpha L_\alpha$:
      Stokes constants within $L_\alpha$ live at slope $-p_\alpha /
      q_\alpha$; cross-edge Stokes constants ($L_\alpha \to L_\beta$ with
      $\alpha < \beta$) are organised along rays of the \emph{higher}
      slope $-p_\beta / q_\beta$.
\end{enumerate}
\end{theorem}

\begin{proof}[Proof sketch]
\textit{Step 1 (Within-edge).} By \Cref{thm:M1}, $\widehat{L} \simeq
\bigoplus L_\alpha$ formally, and the Stokes structure of a direct sum is
the direct sum of per-summand Stokes structures. \Cref{thm:AT1} on the
$q_\alpha$-cover gives $2 p_\alpha \binom{N_\alpha}{2}$ rays; lifting to
the common cover via $s = t^{d/q_\alpha}$ multiplies each ray's
multiplicity by $d / q_\alpha$, yielding $2 P_\alpha \binom{N_\alpha}{2}$
rays at spacing $\pi / P_\alpha$.

\textit{Step 2 (Cross-edge).} On the common cover, $E_\alpha$'s summand
has $Q_i^{(\alpha)}(t) = c_i\, t^{P_\alpha}$ and $E_\beta$'s has
$Q_j^{(\beta)}(t) = c_j\, t^{P_\beta}$. Since slopes are distinct,
$P_\alpha \ne P_\beta$; with $P_\beta > P_\alpha$ and (OS5),
$Q_j^{(\beta)} - Q_i^{(\alpha)} = c_j\, t^{P_\beta} - c_i\, t^{P_\alpha}
+ (\text{lower})$ has leading term $c_j\, t^{P_\beta}$. The anti-Stokes
analysis of \Cref{thm:AT1} applied to this leading term gives $2 P_\beta$
rays per pair.

\textit{Step 3 (Slope filtration).} The Stokes-graded structure of a
slope-filtered irregular connection is standard; see Sabbah
\cite[\S 11]{sabbah-polarizable}, Boalch \cite[\S 3]{boalch-wildhodge}.
\end{proof}

\subsection{\texorpdfstring{$p$-visibility}{p-visibility}}\label{ssec:AT7}

The slope numerator $p$ is invisible to the constellation
(\Cref{rem:N1-analytic}: $\Cl(L)$ depends on $\chi_w$, not on $p$), but
visible to the analytic anti-Stokes skeleton.

\begin{theorem}[$p$-visibility on the cover]\label{thm:AT7}
Fix $q$ and $r$ and let $\chi_w$ be a fixed inner polynomial of degree
$r$. For each $p \in \{p' : 1 \le p' < q,\ \gcd(p', q) = 1\}$, let $L_p$
denote a scalar operator whose Newton polygon has a single dominant edge
of slope $-p/q$ with the prescribed $\chi_w$. Then:
\begin{enumerate}[leftmargin=2em]
\item (Constellation $p$-invariance.) $\Cl(L_p)$ is independent of $p$ as
      a multiset, and the $\Delta$-modulus partition $\{|\Dij|\}$ is
      $p$-independent.
\item ($p$-visibility on the cover.) The anti-Stokes ray set $\Acov_p$
      on the cover $x = t^q$ \emph{depends} on $p$: rays satisfy
      $\arginline \Dij + p\, \phi \equiv \pi/2 \pmod{\pi}$, so varying $p$
      rotates and dilates the ray pattern.
\item (Quantitative formula on the open stratum.) For $\chi_w$ generic
      with $r$ distinct nonzero roots of pairwise distinct moduli and
      $q \ge 3$,
      \[
        \#\Acov_p
        \;=\;
        2 p\, q\, \Big(\, r \;+\; q \tbinom{r}{2}\,\Big).
      \]
      For $q = 2$, $\#\Acov_p = 2 p (r + 4 \binom{r}{2})$.
\end{enumerate}
Hence Phase-III's constellation-level $p$-invariance (P1) holds, while
the analytic anti-Stokes skeleton refines it via the
$p$-proportional ray count.
\end{theorem}

\begin{proof}
\textit{Part (1).} $\chi(c) = \chi_w(c^q)$ depends only on $\chi_w$, so
$\Cl(L_p) = \Roots(\chi)$ is independent of $p$.

\textit{Part (2).} \Cref{thm:AT1} gives $\arginline \Dij + p \phi \equiv
\pi/2 \pmod{\pi}$; solving for $\phi$ shows the spacing $\pi/p$ between
adjacent rays of a fixed pair, manifestly $p$-dependent.

\textit{Part (3) (Counting).} With $\Cl(L) = \bigsqcup_{\alpha=1}^{r}
R_\alpha$, $R_\alpha = \{\zeta_q^j z_\alpha\}$:
\begin{itemize}[leftmargin=2em]
\item \emph{Within-ring.} For each ring, the $\binom{q}{2}$ unordered
      pairs $\{j, k\}$ give $\Dij = z_\alpha\, \zeta_q^k\, (\zeta_q^{j-k}
      - 1)$, whose argument modulo $\pi$ depends only on $(j + k) \mod q$,
      taking $q$ distinct values. Each direction contributes $2p$ rays,
      yielding $2pq$ per ring and $2pqr$ total.
\item \emph{Cross-ring.} The $q^2 \binom{r}{2}$ unordered cross-ring pairs
      generically have pairwise distinct $\arginline \Dij \mod \pi/p$
      (open stratum), each contributing $2p$ rays for $2p \cdot q^2
      \binom{r}{2}$ total.
\end{itemize}
Summing: $\#\Acov_p = 2pqr + 2pq^2 \binom{r}{2} = 2pq(r + q \binom{r}{2})$.
For $q = 2$, the within-ring direction count is $q^2 = 4$ (no overcounting
factor), giving $\#\Acov_p = 2p(r + 4 \binom{r}{2})$.
\end{proof}

\begin{remark}[Phase-IV verification]\label{rem:AT7-verif}
On the PINV family with $q = 5$, $r = 2$, $N = 10$, the formula gives
$\#\Acov_p = 2 p \cdot 5 \cdot (2 + 5) = 70 p$. The Phase-IV measurement
\cite[phase4\_stokes.json]{wkb_dataset_v1} records $70, 140, 210, 280$
for $p = 1, 2, 3, 4$, exact agreement with $70 p$.
\end{remark}

\begin{remark}[Conceptual content]\label{rem:AT7-conceptual}
\Cref{thm:AT7} resolves an apparent tension between P1 (the constellation
does not see $p$) and the analytic Stokes geometry (which manifestly
depends on $p$ via the $t^p$ exponent in $Q_i$). The discrete and
analytic invariants are not in contradiction: the constellation
coarsens away $p$, and the anti-Stokes skeleton recovers it via the ray
spacing.
\end{remark}


%% ============================================================================
%% §8  CONSTELLATION ZOO, TAXONOMY
%% ============================================================================
\section{The 23-class atlas and the constellation zoo}\label{sec:classes}

\subsection{\texorpdfstring{Indexing the image of $\Fclass$}{Indexing the image of F}}\label{ssec:classes-index}

By \Cref{thm:main} and \Cref{cor:ring}, an element of the image of
$\Fclass$ is determined, up to global rotation, by the triple $(q, r,
\pi)$ where $\pi$ is the modulus-with-multiplicity partition of the roots
of $\chi_w$. The slope numerator $p$ is invisible at this level
(\Cref{thm:AT7}(1)), so the atlas is indexed by $(q, r, \pi)$ alone.

\subsection{Taxonomy diagram}\label{ssec:taxonomy}

\begin{figure}[t]
\centering
\fbox{\parbox{0.97\linewidth}{\centering\textsf{[Figure~1
placeholder]}\\[0.4ex]\textsl{Taxonomy diagram of the $23$ constellation
classes; spec in \S\ref{ssec:taxonomy} and the companion dataset
\cite{wkb_dataset_v1}.}}}
\caption{The $\Fclass$-image stratification: $23$ constellation classes.
\textbf{Atlas.} A $4 \times 4$ grid indexed by $q \in \{2, 3, 4\}$
(columns) and $r \in \{1, 2, 3\}$ (rows), with a separate ``$r = 1$''
row spanning all $q$. Each $(q, r)$ cell with $r \ge 2$ is subdivided by
the valuation pattern of $\chi_w$: \emph{simple} (only $\alpha_0,
\alpha_r \ne 0$, single regular $rq$-gon), \emph{double} (one interior
$\alpha_j \ne 0$, two-radius generic or single-radius merger on the
discriminant locus), \emph{mixed/generic} (all $\alpha_j \ne 0$).
\textbf{Generic vs.\ special encoding.} Solid borders mark the generic
stratum, dashed grey borders the real-coefficient stratum, grey fill the
single-ring collapses, corner badges the accidental-symmetry classes.
\textbf{Degeneracy island.} A detached $1 \times 4$ inset records the
four stress-test classes (\texttt{STRESS\_equalmod},
\texttt{STRESS\_doubleroot}, \texttt{STRESS\_zeroroot},
\texttt{STRESS\_accidental\_sym}) with codimension-$1$ connector arrows
to the parent atlas cells.}
\label{fig:taxonomy}
\end{figure}

The taxonomy diagram of \Cref{fig:taxonomy} encodes the image
stratification visually. A precise textual specification (axes, sub-tile
structure, encoding conventions, degeneracy island, legend) is recorded
in the dataset \cite[\texttt{taxonomy.txt}]{wkb_dataset_v1} and is
suitable for a figure-rendering pipeline. The diagram lives in this
section in the camera-ready manuscript.

\subsection{Constellation zoo}\label{ssec:zoo}

\begin{figure}[t]
\centering
\includegraphics[width=0.95\linewidth]{figures/constellations.png}
\caption{Representative constellations across the $\Fclass$-image. The
$3 \times 4$ grid covers the principal strata. Top row (A1--A4): simple /
binomial classes --- single regular polygon (a single Galois orbit of
$\chi_w$). Middle row (B1--B4): double and generic-complex classes ---
two or three concentric $\mu_q$-orbits. Bottom row (C1--C4): the four
named degenerate strata. Panels A3 and C4 are visually identical $6$-gons
but live in different cells: A3 at slope $1/3$ has minimum symmetry $C_3$
enhanced to $C_6$; C4 at slope $1/2$ has $C_2$ enhanced to $C_6$. The
constellation alone does not recover the cell --- the Newton-polygon
denominator $q$ is required. This is the visible signature that
$\Fclass$ is many-to-one and that accidental symmetry is a feature of
every simple-class operator. Source:
\cite[\texttt{constellations.png}, \texttt{plot\_constellations.py}]{wkb_dataset_v1}.}
\label{fig:zoo}
\end{figure}

\Cref{fig:zoo} shows twelve representative panels. The full enumeration
of all $23$ classes, with explicit operator representatives, $\chi_w$
data, and panel pointers, lives in the companion dataset
\cite[\texttt{clusters.json},
\texttt{constellations.json}]{wkb_dataset_v1}.


%% ============================================================================
%% §9  RELATED WORK
%% ============================================================================
\section{Related work}\label{sec:related}

\paragraph{ADE singularities and discrete classification.}
As a methodological precedent we take the ADE classification of simple
singularities (Arnol'd \cite{arnold1976}, Brieskorn \cite{brieskorn}) and
of finite subgroups of $SU(2)$ (McKay \cite{mckay}): a small discrete
invariant---a Dynkin diagram---organises an apparently heterogeneous zoo
of analytic objects, with strata governed by simple-root combinatorics
and degenerations corresponding to root-system collisions. Our $23$
constellation classes occupy an analogous slot: a finite combinatorial
atlas indexed by the discrete data $(q, r, \pi, \chi_w)$, with degenerate
strata produced by collisions among the inner roots $w_i$ (the analogue
of root coincidences). Unlike the ADE setting, our classes are not
labelled by Dynkin diagrams---the relevant combinatorics is the partition
of $\chi_w$'s roots by modulus and Galois orbit---but the \emph{style} of
classification (discrete invariants, stratified moduli, degeneration
loci) is closely modelled on it.

\paragraph{Formal classification: Levelt--Turrittin and descendants.}
The Levelt--Turrittin theorem (Levelt \cite{levelt}, Turrittin
\cite{turrittin}; modern accounts in Sabbah \cite{sabbah-isomon} and
Sabbah--Mochizuki \cite{mochizuki-asymptotic}) gives a formal
classification of meromorphic connections by \emph{irregular type}: a
multiset of exponential factors $e^{P_i(x^{1/q})}$ together with
regular-singular twists. Birkhoff--Trjitzinsky \cite{birkhoff-trjitzinsky}
and the $q$-difference theories of Sauloy \cite{sauloy}, Ramis
\cite{ramis-stokes,ramis-q-stokes}, and Zhang \cite{zhang} play the
parallel role for difference and $q$-difference operators. Each of these
theories classifies \emph{the operator} up to formal or analytic
equivalence; our object is different. We classify the \emph{image} of
the WKB-extraction functor $\Fclass$, \ie\ the discrete geometry that
$\NL$ and its leading-edge coefficients deposit on the complex plane.
The connection is one-way: the Levelt--Turrittin/Birkhoff--Trjitzinsky
exponents are recovered from $\Cl(L)$ together with auxiliary phase data,
but $\Cl(L)$ itself is coarser and admits a clean independent
classification.

\paragraph{Stokes graphs and spectral networks.}
A third neighbour is the Stokes-graph classification of Painlev{\'e} and
other low-rank equations (Kawai--Takei \cite{kawai-takei},
Iwaki--Nakanishi \cite{iwaki-nakanishi}, Aoki--Kawai--Takei
\cite{aoki-kawai-takei}), and more broadly the spectral-network programme
of Gaiotto--Moore--Neitzke \cite{gmn-spectral}. There, classification
proceeds via the \emph{combinatorial topology of Stokes lines} in the
complex plane, which depends on $c_i - c_j$ rather than $|c_i|$. Our
ring-modulus partition is therefore \emph{not} a Stokes invariant and
does not reproduce spectral-network combinatorics; instead it captures
the strictly weaker but exact information of the multiset $\{c_i^q\}$.
This makes $\Fclass$ a coarsening of the spectral-network functor, which
we view as an asset: the coarser invariant is computable directly from
Newton-polygon data and admits a closed-form classification, whereas
Stokes graphs in general do not.

\paragraph{Wild character varieties and irregular Riemann--Hilbert.}
The geometry of \emph{wild character varieties} and the irregular
Riemann--Hilbert correspondence (Boalch
\cite{boalch-stokes,boalch-wildhodge}, Sabbah \cite{sabbah-polarizable},
Mochizuki \cite{mochizuki-asymptotic}) produces moduli stacks parametrising
local data of irregular connections; our $\Fclass$-image is a discrete
shadow of such a stack restricted to scalar operators with one dominant
edge. We do not attempt to lift the present classification to the level
of stacks or to incorporate Stokes structures intrinsically; rather, we
exhibit a self-contained classification of the discrete part of the local
data, with the analytic refinement of \S\ref{sec:analytic} attaching the
first layer of Stokes information (anti-Stokes directions; slope-graded
Stokes filtration) but stopping short of explicit Stokes constants. The
contribution is best understood as building a separately classifiable
layer of the local picture---a Newton-polygon WKB stratum---on which the
fuller analytic theories can subsequently be erected.


%% ============================================================================
%% §10  DISCUSSION AND VISION
%% ============================================================================
\section{Discussion and future directions}\label{sec:vision}

The natural object of study, in our view, is not any individual operator
but the constellation functor $\Fclass$ itself: the map from
locally-irregular linear models to their WKB exponent multisets. Once
$\Fclass$ is named, the questions become geometric---what is its image,
what are its fibres, what are its strata, how does it deform---and admit
answers in closed form when restricted to one dominant Newton edge. The
$23$ classes catalogued here populate the lowest-dimensional cells of
$\mathrm{Im}(\Fclass)$; their adjacencies trace out the discriminant,
equal-modulus, and zero-root walls along which the structure changes.

The natural sequel is to walk the same programme up the complexity
ladder. \Cref{thm:M1,thm:N1,thm:AT5} already step one rung: multi-edge
Newton polygons (where $\Fclass$ acquires a slope-filtration target) and
nested ramification (where the visible symmetry inflates from $\mu_q$ to
$\mu_{\mathrm{lcm}(q, q')}$). Beyond:
\begin{itemize}[leftmargin=2em]
\item \textbf{Coupled systems.} Matrix and meromorphic-connection systems
      have root systems on each edge; their WKB target is a wild
      character variety. \Cref{thm:S1} treats only the block-diagonal
      case. Genuinely coupled systems (non-semisimple leading symbol)
      lie outside the present classification and connect naturally to
      Boalch's wild non-abelian Hodge theory \cite{boalch-wildhodge}.
\item \textbf{Stokes constants and the analytic refinement.} The
      anti-Stokes skeleton of \Cref{thm:AT1,thm:AT5,thm:AT7} attaches
      \emph{directions} to constellations; the next layer is the
      attachment of \emph{Stokes constants} along each ray. The locus of
      Stokes-equivalent operators above a fixed constellation is the
      kernel of a refined functor $\widetilde\Fclass$ whose image is a
      wild-character-variety stratum.
\item \textbf{Difference and $q$-difference Stokes lifts.} The analytic
      theorems are stated for $\op = \opD$; the analogues for $\opS$
      (Birkhoff--Trjitzinsky) and $\opT$ (Ramis $\theta$-Stokes) require
      separate analytic foundations and are deferred. The constellation
      identities (\Cref{thm:main,thm:M1,thm:N1,thm:D2,thm:S1}) hold
      verbatim across equation types.
\item \textbf{Composite-$q$ Kummer dichotomy.} \Cref{thm:D2} is sharp for
      prime $q$ within the Phase-III search window. The composite-$q$
      version requires the refinement of \Cref{rem:D2-composite} and
      the transcendental-degree statement of \Cref{rem:D2-search}.
\end{itemize}

We expect the functorial reframing to be useful precisely because it
makes a single, universal target geometry visible behind a great many
local stories that, treated one at a time, have looked unrelated. The
classifiable object is the functor; the discrete combinatorics of the
constellation sits transversally to the analytic theory---coarser,
computable, and structurally separable.

%% ============================================================================
%% §11  REPRODUCIBILITY NOTE
%% ============================================================================
\section{Reproducibility note}\label{sec:reproducibility}

All computational artifacts---the operator atlas, characteristic
polynomials, constellations, clusters, conjectures, counterexample
catalogue, taxonomy specification, and rendered figures---are bundled in
the companion Zenodo dataset \cite{wkb_dataset_v1}. The dataset is
organised under the directory structure
\begin{center}
\texttt{/code\ /data\ /figures\ /reports\ /theorems\ /paper-draft\ /fleets}
\end{center}
with \texttt{README.md} at the root. The five fleet pipelines
(\texttt{fleet.py}, \texttt{fleet3.py}, \texttt{fleet4.py},
\texttt{fleet5.py}) are deterministic and self-contained
(\texttt{Python 3.10+}, \texttt{sympy}, \texttt{numpy},
\texttt{matplotlib}); no random seeds are used. Running each fleet in
sequence reproduces the JSON atlas and the rendered constellation zoo
(\texttt{constellations.png}).

The proofs of \Cref{thm:M1,thm:N1,thm:D2,thm:S1,thm:AT1,thm:AT3,thm:AT5,thm:AT7}
are also bundled in the dataset as the file \texttt{proofs.md}, in
long-form expository style with explicit cross-references to
Phase-III/IV evidence. The integration map
(\texttt{integration\_map.md}) records, for each theorem, the dataset
artifact that supplies its empirical evidence. The counterexample
catalogue (\texttt{counterexamples.json}) records the strata on which
naive forms of the theorems fail; each theorem statement above includes
a remark guarding against the relevant counterexamples.

\medskip
\noindent\textbf{Citation.} If you use this work, please cite both the
paper and the dataset; the dataset DOI (placeholder
\texttt{10.5281/zenodo.XXXXXXX}) will be assigned by Zenodo on
publication of v1.0.

%% ============================================================================
%% COMPUTATIONAL AND AI DISCLOSURE
%% ============================================================================
\section*{Computational and AI Disclosure}\label{sec:disclosure}

\paragraph{Computational environment.}
All numerical and symbolic work in this paper was carried out in
\texttt{Python 3.10+} using \texttt{sympy} (exact symbolic algebra),
\texttt{fractions} and \texttt{numpy} (exact rational and floating-point
arithmetic), and \texttt{matplotlib} (figure rendering). No randomised
algorithms were used and no random seeds were set: every JSON artifact
in the companion dataset \cite{wkb_dataset_v1} is reproduced bytewise on
re-execution. The manuscript was typeset with MiKTeX/\TeX{} Live using
the package set listed in the preamble; the bibliography was processed
with \texttt{bibtex} against \texttt{references.bib}.

\paragraph{Automated pipelines (Phases I--V).}
The computational atlas was produced by five fleet scripts, each
implementing one phase of the investigation:
\begin{itemize}[leftmargin=2em,itemsep=2pt,topsep=2pt]
\item \texttt{fleet.py} (Phases I/II) --- generates the single-edge
  operator atlas, computes Newton polygons, extracts characteristic
  polynomials $\chi$ and inner polynomials $\chi_w$, clusters the
  exponent constellations, and emits the initial conjecture list.
\item \texttt{fleet3.py} (Phase III) --- extends the atlas to
  multi-edge polygons, $2\times 2$ and $3\times 3$ systems, resonant
  slopes ($q\le 7$), and nested ramification; computes
  $\Delta_{ij}$ partitions, Galois-orbit invariants, and the
  inner-polynomial $\mathbb{Z}$-resonance dichotomy.
\item \texttt{fleet4.py} (Phase IV) --- lifts the discrete invariants
  to the analytic layer via the Levelt--Turrittin decomposition,
  predicts anti-Stokes rays from the $\Delta_{ij}$ geometry, and
  cross-checks the predictions against numerical Stokes computations.
\item \texttt{fleet5.py} (Phase V) --- consolidates the conjectures
  into theorem--proof pairs, runs each statement against the
  counterexample catalogue, and emits the proof skeletons recorded in
  \texttt{proofs.md} with placements specified by
  \texttt{integration\_map.md}.
\end{itemize}
The pipelines are deterministic and self-contained; running them in
sequence reproduces every JSON file, the rendered constellation zoo
(\texttt{constellations.png}), and the proof skeletons.

\paragraph{Role of AI assistance.}
This work was conducted as independent research with AI-assisted
execution (GitHub Copilot CLI driving an Anthropic Claude Opus~4.7
model). The assistant was used in five clearly demarcated roles:
\begin{enumerate}[label=(\roman*),leftmargin=2em,itemsep=2pt,topsep=2pt]
\item \emph{Conjecture selection.} The assistant scanned the
  Phase-I--III invariant tables for candidate universality laws; the
  author selected which conjectures to promote to the main text.
\item \emph{Proof-skeleton generation.} For each promoted conjecture
  the assistant produced a draft proof skeleton naming the algebraic
  ingredients (Newton-polygon combinatorics, $\mu_q$-pullback,
  $\chi_w$ factorisation, Levelt--Turrittin, determinant polygon);
  the author verified each step and rewrote the load-bearing
  arguments.
\item \emph{Counterexample guarding.} The assistant cross-referenced
  each candidate theorem against \texttt{counterexamples.json} and
  proposed genericity hypotheses (prime $q$, single-edge, non-resonant
  inner polynomial) that survive the catalogued exceptions; the
  author adjudicated which restrictions were accepted into the final
  statements.
\item \emph{Refinement passes.} The assistant unified notation, fixed
  cross-references, and polished prose; all theorem statements,
  remarks, and limitations were reviewed by the author.
\item \emph{Manuscript assembly.} The assistant inserted the
  theorem--proof blocks into the section placements recorded in
  \texttt{integration\_map.md}, assembled the bibliography, and
  produced the arXiv and journal submission bundles.
\end{enumerate}
No proprietary or closed-source models were used beyond the
publicly-available AI assistant named above; no model was fine-tuned on
the project's data, and no AI-mediated literature summaries were used
in place of primary sources for citations.

\paragraph{Human accountability.}
All mathematical statements, theorem formulations, proofs, and
classification claims in this paper were reviewed, validated, and
finalised by the author. In particular, the structural law of
\Cref{thm:main}, the multi-edge and nested-ramification factorisations
of \Cref{thm:M1,thm:N1}, the prime-$q$ Kummer dichotomy of
\Cref{thm:D2}, the system pullback of \Cref{thm:S1}, the four
analytic-lifting theorems \Cref{thm:AT1,thm:AT3,thm:AT5,thm:AT7}, and
the 23-class taxonomy of \Cref{sec:classes} reflect human mathematical
judgment. The author retains sole accountability for all claims,
methodology choices, and the published result.

\paragraph{Reproducibility.}
The complete computational record --- operators, Newton polygons,
characteristic polynomials, constellations, clusters, counterexamples,
Phase-III/IV invariants, proof skeletons, integration map, and
rendered figures --- is deposited as the companion Zenodo dataset
\cite{wkb_dataset_v1}, licensed CC-BY-4.0 (text and data) and MIT
(code and pipelines). The dataset DOI is \texttt{10.5281/zenodo.XXXXXXX}
(placeholder, assigned by Zenodo at publication of v1.0). See
\Cref{sec:reproducibility} for the command sequence that reproduces
every artifact.


%% ============================================================================
%% ACKNOWLEDGEMENTS
%% ============================================================================
\section*{Acknowledgements}\label{sec:ack}

The fleet pipeline that generated the computational atlas was prototyped
in successive phases of automated investigation; the author thanks the
contributors to \texttt{sympy}, \texttt{numpy}, and \texttt{matplotlib},
on which the pipeline rests. Acknowledgements of individual collaborators
and funding sources will be added in the camera-ready version (v2.0).


%% ============================================================================
%% BIBLIOGRAPHY
%% ============================================================================
\bibliographystyle{alpha}
\bibliography{references}

\end{document}
